\chapter{Rejestr stanu książki}
\markboth{Rejestr stanu książki}{Rejestr stanu książki}
Wszystkie liczby w tym dodatku są \emph{liczone} z danych i źródeł (\code{tools/gen.py}), nie wpisywane. Ten rejestr istnieje po to, by czytelnik wiedział,
czego książka jeszcze nie ma.

\section*{Zawartość}
\begin{center}\small
\begin{tabular}{@{}lrrrrr@{}}
\toprule
 & \textsf{\textbf{A1}} & \textsf{\textbf{A2}} & \textsf{\textbf{B1}} & \textsf{\textbf{B2}} & \textsf{\textbf{razem}}\\
\midrule
jednostki & \val{units.A1} & \val{units.A2} & \val{units.B1} & \val{units.B2} & \val{units.total}\\
napisane jednostki & \val{written.A1} & \val{written.A2} & \val{written.B1} & \val{written.B2} & \val{written.total}\\
słownictwo (wpisy) & \val{vocab.A1} & \val{vocab.A2} & \val{vocab.B1} & \val{vocab.B2} & \val{vocab.total}\\
ćwiczenia (razem z generowanym treningiem) & \val{exercises.A1} & \val{exercises.A2} & \val{exercises.B1} & \val{exercises.B2} & \val{exercises.total}\\
w tym trening słownictwa (generowany) & \val{drills.A1} & \val{drills.A2} & \val{drills.B1} & \val{drills.B2} & \val{drills.total}\\
dialogi & \val{dialogs.A1} & \val{dialogs.A2} & \val{dialogs.B1} & \val{dialogs.B2} & \val{dialogs.total}\\
cele can-do & \val{candos.A1} & \val{candos.A2} & \val{candos.B1} & \val{candos.B2} & \val{candos.total}\\
ramki \emph{Pułapka} & \val{traps.A1} & \val{traps.A2} & \val{traps.B1} & \val{traps.B2} & \val{traps.total}\\
ramki \emph{Szwajcaria} & \val{swiss.A1} & \val{swiss.A2} & \val{swiss.B1} & \val{swiss.B2} & \val{swiss.total}\\
ramki \emph{W pracy} & \val{work.A1} & \val{work.A2} & \val{work.B1} & \val{work.B2} & \val{work.total}\\
\bottomrule
\end{tabular}
\end{center}

Fiszki Beamer (\code{flashcards/}): \val{cards.grammar.total} kart gramatycznych (pisanych ręcznie, A1 \val{cards.grammar.A1}, A2 \val{cards.grammar.A2},
B1 \val{cards.grammar.B1}, B2 \val{cards.grammar.B2}) i \val{cards.vocab.total} kart słownictwa (generowanych z \code{data/vocab/}); fiszki Anki mają te same słowa w obie strony.

Rodzaje rzeczowników w słowniku: \artm{der} \val{nouns.der}, \artf{die} \val{nouns.die}, \artn{das} \val{nouns.das}; czasowników \val{vocab.verbs}.
Tygodni w planie (z buforowymi): \val{weeks.total}.

\section*{Dług, który książka ujawnia}
\begin{itemize}
\item \textbf{Weryfikacja przez native speakera: NIE ZOSTAŁA PRZEPROWADZONA.} Niemieckie przykłady, dialogi i klucze odpowiedzi nie były czytane przez
rodzimego użytkownika języka. To jedyny rejestr, który nie jest licznikiem, tylko deklaracją, i nie może zniknąć po cichu.
\item \textbf{Fakty wymagające źródła pierwotnego: \val{verify.total}} miejsc oznaczonych \verify{w ten sposób}. Dotyczą prawa pracy, urzędów, podatków i
formatów egzaminów w Szwajcarii. Każde trzeba sprawdzić u źródła przed decyzją.
\item \textbf{Słownictwo.} \val{vocab.total} wpisów to rdzeń, nie pełny zasób B2; uzupełniaj z tekstów (Dodatek~E).
\item \textbf{Brak nagrań audio.} Książka nie zawiera dźwięku. Wymowę ćwicz z DW, SRF i korepetytorem.
\item \textbf{Czas.} Plan zakłada ok.\ 600 godzin w roku; przy 60 minutach dziennie uzyskasz poziom B1 (Wstęp).
\end{itemize}
